\begin{table}[H]
\centering
\footnotesize
\caption{Countries where observed employment composition most raises or lowers the occupation-layer exposed share.}
\resizebox{\textwidth}{!}{%
\begin{tabular}{p{0.12\textwidth}p{0.25\textwidth}p{0.19\textwidth}ccc}
\toprule
Direction & Country & Income group & Linkage-weighted & Employment-weighted & Adjustment \\
\midrule
Raised & Japan & High income & 0.60 & 0.64 & +0.04 \\
 & Korea, Rep. & High income & 0.60 & 0.64 & +0.04 \\
 & Hong Kong SAR, China & High income & 0.56 & 0.59 & +0.03 \\
 & Singapore & High income & 0.61 & 0.63 & +0.03 \\
 & Italy & High income & 0.55 & 0.57 & +0.02 \\
 & United Kingdom & High income & 0.57 & 0.59 & +0.02 \\
Lowered & Rwanda & Low income & 0.30 & 0.22 & -0.09 \\
 & Benin & Lower middle income & 0.25 & 0.16 & -0.08 \\
 & Mozambique & Low income & 0.19 & 0.11 & -0.08 \\
 & Uganda & Low income & 0.21 & 0.13 & -0.08 \\
 & Cote d'Ivoire & Lower middle income & 0.29 & 0.21 & -0.08 \\
 & Zambia & Lower middle income & 0.26 & 0.19 & -0.08 \\
\bottomrule
\end{tabular}
}
\caption*{\scriptsize Note: Positive adjustments mean that more exposed ISCO major groups account for a larger employment share than under the task-to-occupation linkage weights in that country.}
\label{tab:appendix_ilostat_country_adjustment}
\end{table}
